% page 572 of the facsimile (image p-571 of the scan)
% block 157.5 x 237.7 mm ; left margin 8.0 mm
\pgc{-12.700mm}{0.000mm}{
\l{\cel{3}564}
\l{}
\l{\sou{surtut}o.\cel{2}Vesto qua konservas la kaloro di la korpo, quan}
\l{\cel{3}la viri +weras sur sua frako, redingoto. - EFIS.}
\l{}
\l{\sou{surveyar}. (\sou{trans.)}\cel{2}Sorgar per regardar kun atenco sen-inter-}
\l{\cel{3}rupta, time di defekto, di eroro, di neglijo, e c.- EFI.}
\l{}
\l{\sou{sus}\cel{1}!\cel{2}Interjeciono qua uzesas por ecitar ad inter-defenso,}
\l{\cel{3}a persequo.}
\l{}
\l{\sou{suskriptar}. (\sou{trans., pri}) Engajar su pri pagor parto (di}
\l{\cel{3}ulo). - DEFIS.}
\l{}
\l{\sou{suspektar}. (\sou{trans., pri}).I. (a) Konjektar, de ula aspekti, ke}
\l{\cel{3}(ulu) esas kulpozo ; (b) Admisar la existo, pro ula aspek-}
\l{\cel{3}ti, (di lo mala pri qua onu konjektas ke ulu esas la au-}
\l{\cel{3}toro). - II. Konjektar (ulo) de ula aspekti. - III. (\sou{metaf}.}
\l{\cel{3}Dicernar instinte. - EFIS.}
\l{}
\l{\sou{suspendar}. (\sou{trans.}) Tenar super sulo (persono, kozo) tale ke}
\l{\cel{3}lu pendas. - DEFIS.}
\l{}
\l{\sou{suspensar}. (\sou{trans.)}\cel{2}Revokar, dure di nur kelka tempo,de lua}
\l{\cel{3}oficiereso. - FI}
\l{}
\l{\sou{suspensoro}. (\sou{biol.}) Ligamento qua tenas suspendita. -}
\l{}
\l{\sou{suspensorio}. Bandajo qua sustenas la testikuli. - DEFIS.}
\l{}
\l{\sou{sustenar}. (\sou{trans.}) Subtenar (ulo, ulu) per portar parto de}
\l{\cel{3}lua pezo. - EFIS.}
\l{}
\l{\sou{sustracionar}. (\sou{trans.})\cel{2}(\sou{aritm.}) Deprenar nombro de altra,}
\l{\cel{3}por determinar lia interdifero.- DEFIS.}
\l{}
\l{\sou{sustraktar}. (\sou{trans.,}\cel{1}de)\cel{2}I. Deprenar (ulo) de ulu, por ke lu}
\l{\cel{3}ne povez uzar olu. - II. Deprenar (ulo qua apartenas ad}
\l{\cel{3}ulu) cele, per habileso o fraudo. - FIS.}
\l{}
\l{\sou{susurar}. (\sou{trans., ulo, ad)}\cel{1}Parolar deslaute, preske an la}
\l{\cel{3}orelo (di ulu). -FI.}
\l{}
\l{\sou{sutar}. (\sou{trans.}) Asemblar, ligar, per filo, kordeto, quan onu}
\l{\cel{3}igas trapasar per agulo, per puncilo. - EFI.}
\l{}
\l{\sou{sutano}. Vesto butonagita de supre ad infre, qua kovras til la}
\l{\cel{3}pedi, quan olim +weris la mediki, la sacerdoti, e quan nur}
\l{\cel{3}ici duras +werar. - DEFIS.}
\l{}
\l{\sou{suverena}. I. Qua havas la autoritato suprega en la stato.-}
\l{\cel{3}II. (\sou{politiko-yuro}) La\cel{1}\sou{suvereno}\cel{1}: La persono (personi) en}
\l{\cel{3}rua (qui) rezidas la autoritato suprega leg-exekutigala}
\l{\cel{3}(la rejo, en monarkio ; la populo, en demokratio ; la}
\l{\cel{3}grandi, en oligarkio). - DEFIS.}
}
